% Generated draft: requires mathematical and source review.
\documentclass[11pt]{article}
\usepackage[margin=25mm]{geometry}
\usepackage{fontspec}
\setmainfont{Latin Modern Roman}
\usepackage{amsmath,amssymb,mathtools,unicode-math}
\setmathfont{Latin Modern Math}
\usepackage{enumitem,needspace,xurl,hyperref}
\setlength{\parindent}{0pt}
\setlength{\parskip}{5pt}
\setlength{\emergencystretch}{3em}
\newcommand{\N}{\mathbb{N}}
\newcommand{\Z}{\mathbb{Z}}
\newcommand{\Q}{\mathbb{Q}}
\newcommand{\R}{\mathbb{R}}
\newcommand{\C}{\mathbb{C}}
\newcommand{\F}{\mathbb{F}}
\newcommand{\A}{\mathbb{A}}
\newcommand{\PP}{\mathbb{P}}
\newcommand{\E}{\mathbb{E}}
\newcommand{\eps}{\varepsilon}
\newcommand{\dd}{\,\mathrm{d}}
\begin{document}
\section*{UnsolvedMath ID: 700794}
\textbf{Source problem code:} catalogue entry
\subsection{Definitions and mathematical statement}
For a group \(G\), define its derived series by \(G^{(0)}=G\) and \(G^{(n+1)}=[G^{(n)},G^{(n)}]\) for each integer \(n\geq 0\), where \([H,H]\) denotes the subgroup generated by the commutators of elements of \(H\). A group \(G\) is \emph{soluble} if \(G^{(n)}=\{1\}\) for some integer \(n\geq 0\), and \emph{insoluble} otherwise. It is \emph{residually soluble} if, for every \(g\in G\setminus\{1\}\), there exists a normal subgroup \(N\trianglelefteq G\) such that \(g\notin N\) and \(G/N\) is soluble. These soluble quotients need not be finite, and their derived lengths need not have a common bound. The \emph{maximum condition on subgroups} means that every ascending chain
\[
H_1\leq H_2\leq H_3\leq\cdots\leq G
\]
stabilizes: there is an integer \(m\geq 1\) such that \(H_n=H_m\) for all \(n\geq m\). This condition concerns all subgroups, not merely normal subgroups.

\textbf{Question (J. Wiegold, Kourovka Notebook 11.102).} Does there exist a group \(G\) which is residually soluble, satisfies the maximum condition on subgroups, and satisfies
\[
G^{(n)}\neq\{1\}\qquad\text{for every integer }n\geq 0?
\]
\subsection{Short English statement}
Does there exist an insoluble group whose nonidentity elements can each be detected in a soluble quotient and in which every ascending chain of subgroups stabilizes? The supplied registry labels this problem open; that label is not independently verified here.
\subsection{Sources}
\begin{itemize}
\item Supplied local UnsolvedMath registry record, permanent ID \(700794\), exact source code \texttt{catalogue entry}, titled \emph{Kourovka Notebook Problem 11.102}. The record was updated on \texttt{2026-10-16T14:00:00Z} and carries the status \texttt{open}; this is supplied registry metadata, not independent verification.
\item \emph{The Kourovka Notebook (Unsolved Problems in Group Theory)}, original problem 11.102, proposed by J. Wiegold in the 11th issue (1990). The supplied source identifies the 21st issue edition (2026), arXiv:1401.0300v48, and describes the collection as edited by E. I. Khukhro and V. D. Mazurov. This bibliographic attribution is reproduced from the supplied record; the referenced edition has not been independently consulted.
\end{itemize}
\end{document}
